\scanpage{32}
\mspage{12}

following properties (characterizing a \uline{strict} Quillen triple)

4) \deleted{\unsure{A|At} any base change} \inserted{$W$ stable by base changes} by any $g\in F$, and by cobase changes by any $g\in C$.

\deleted{\unsure{g}}

\uline{NB} From my point of view, $W$ is given beforehand in a rather tangible way, by cohomological criteria, which turn out very handy. In Kan's-Quillen's point of view, $W$ was a little more hidden maybe in terms of (say) its structure, they get to it via \deleted{\unsure{a}} the pair $(C,TF)$ and $F\supset TF$ the Kan fibrations, which is a somewhat ad-hoc construction in terms of a certain \deleted{\unsure{g}} a-priori generating set $(TC)_0$ --- corresponding roughly to expressing the \unsure{horn} lifting property. What remains to be understood, in the context say of general test categories, is the relationship between this and the direct description $TC=C\cap W$, \deleted{\unsure{what}} i.e why this is just the \uline{saturation} of $(TC)_0$. Maybe I should have a look at Gabriel-Zisman and their treatment of ``anodyne extensions''.

\mspage{31}

